\section{习题}

\subsection{本章速览}

\begin{itemize}
\item \textbf{有解判定：}$Ax=b$ 有解 $\iff r(A)=r(A\mid b)$。有解时，$r(A)=n$ 为唯一解，$r(A)<n$ 为无穷多解；若 $r(A\mid b)=r(A)+1$ 则无解。
\item \textbf{齐次方程组：}$Ax=0$ 恒有零解；有非零解 $\iff r(A)<n$。解集是 $\mathbb R^n$ 的子空间，维数为 $n-r(A)$。
\item \textbf{基础解系：}它是解空间的一组基，恰含 $n-r(A)$ 个向量。把系数矩阵化为行最简形，令自由未知量依次取单位向量，即可读出一组基础解系。
\item \textbf{非齐次通解结构：}$x=\eta^*+k_1\xi_1+\cdots+k_{n-r}\xi_{n-r}$，即“一个特解加导出组的通解”。两个解之差是导出组的解；解集是仿射子空间 $\eta^*+\ker A$，当 $b\neq0$ 时不是向量空间。
\item \textbf{含参讨论：}方阵先算 $\det A$ 找分界点，再对每个分界值把增广矩阵化为行阶梯形定秩；非方阵只能直接对增广矩阵消元。判断有解必须比较 $r(A)$ 与 $r(A\mid b)$，不能只看 $\det A$。
\item \textbf{解的结构上限：}解集中线性无关的向量最多 $n-r+1$ 个；$Ax=b$ 有唯一解 $\iff r(A)=r(A\mid b)=n$。
\end{itemize}

\subsection{基础过关题}

\textbf{1.} 设 $A$ 为 $m\times n$ 矩阵，$b$ 为 $m$ 维列向量。判断下列说法是否正确，并说明理由：
（1）若 $Ax=b$ 有两个不同的解，则它必有无穷多解；
（2）若 $Ax=0$ 只有零解，则 $Ax=b$ 有唯一解；
（3）若 $r(A)=r(A\mid b)$，则 $Ax=b$ 有解；
（4）若 $m<n$，则 $Ax=0$ 必有非零解。

\ifshowsolutions%
\boxed{\text{解}}

\textbf{思路：}逐条回到秩判定与解的性质。命题（2）最容易误判：齐次组只有零解只说明列满秩，并不能保证非齐次组有解。

（1）\textbf{正确。}设 $x_1\neq x_2$ 都是 $Ax=b$ 的解，则对任意实数 $t$，
\[
A\bigl(x_1+t(x_2-x_1)\bigr)=Ax_1+t(Ax_2-Ax_1)=b+t(b-b)=b,
\]
即 $x_1+t(x_2-x_1)$ 仍是解。由于 $x_2-x_1\neq0$，不同的 $t$ 给出不同的解，故有无穷多解。

（2）\textbf{错误。}$Ax=0$ 只有零解只说明 $r(A)=n$（列满秩），而非齐次组可能无解。取
\[
A=\begin{pmatrix}1\\0\end{pmatrix},\qquad b=\begin{pmatrix}0\\1\end{pmatrix},
\]
则 $Ax=0$ 即 $x=0$，只有零解；但 $Ax=b$ 要求同时满足 $x=0$ 与 $0=1$，故无解。

（3）\textbf{正确。}这正是非齐次线性方程组有解的充要条件（Kronecker--Capelli 定理）。

（4）\textbf{正确。}因为 $r(A)\le\min\{m,n\}=m<n$，所以解空间维数 $n-r(A)>0$，齐次组除零解外还有非零解。

\boxed{\text{答案：（1）正确；（2）错误；（3）正确；（4）正确。}}
\else \answerarea[8] \fi%

\textbf{2.} 求齐次线性方程组
\[
\begin{cases}
x_1+x_2+x_3+x_4=0,\\
x_1+2x_2+3x_3+4x_4=0,\\
2x_1+3x_2+4x_3+5x_4=0
\end{cases}
\]
的一个基础解系和通解。

\ifshowsolutions%
\boxed{\text{解}}

\textbf{思路：}对系数矩阵作初等行变换化为行最简形，主元变量用自由未知量表示，再令自由未知量取单位向量读出基础解系。

\[
A=\begin{pmatrix}
1&1&1&1\\
1&2&3&4\\
2&3&4&5
\end{pmatrix}
\xrightarrow[R_3-2R_1]{R_2-R_1}
\begin{pmatrix}
1&1&1&1\\
0&1&2&3\\
0&1&2&3
\end{pmatrix}
\xrightarrow{R_3-R_2}
\begin{pmatrix}
1&1&1&1\\
0&1&2&3\\
0&0&0&0
\end{pmatrix}.
\]
再作 $R_1-R_2$：
\[
\xrightarrow{R_1-R_2}
\begin{pmatrix}
1&0&-1&-2\\
0&1&2&3\\
0&0&0&0
\end{pmatrix}.
\]
于是 $r(A)=2$，未知数个数 $n=4$，基础解系含 $n-r=2$ 个向量，$x_3,x_4$ 为自由未知量。由行最简形得
\[
x_1=x_3+2x_4,\qquad x_2=-2x_3-3x_4.
\]
令 $(x_3,x_4)=(1,0)$，得 $\xi_1=(1,-2,1,0)^T$；令 $(x_3,x_4)=(0,1)$，得 $\xi_2=(2,-3,0,1)^T$。

直接代回可验证：$\xi_1,\xi_2$ 都满足三个方程，且二者线性无关（第 3、4 分量构成的子式不为零），个数又恰为 $n-r=2$，故它们构成一个基础解系。

通解为
\[
x=k_1\xi_1+k_2\xi_2,\qquad k_1,k_2\in\mathbb R.
\]
基础解系不唯一，例如把 $\xi_1$ 乘以任意非零常数仍可与其他向量组成基础解系，但向量个数恒为 $2$。

\boxed{\text{答案：}\xi_1=(1,-2,1,0)^T,\ \xi_2=(2,-3,0,1)^T;\ x=k_1\xi_1+k_2\xi_2}
\else \answerarea[10] \fi%

\textbf{3.} 求非齐次线性方程组
\[
\begin{cases}
x_1+x_2+x_3=3,\\
2x_1-x_2+x_3=2,\\
3x_1+2x_3=5
\end{cases}
\]
的通解。

\ifshowsolutions%
\boxed{\text{解}}

\textbf{思路：}先对增广矩阵消元定秩，确认有解后取一个自由未知量，写出特解与导出组的基础解系。

\[
(A\mid b)=\left(\begin{array}{ccc|c}
1&1&1&3\\
2&-1&1&2\\
3&0&2&5
\end{array}\right)
\xrightarrow[R_3-3R_1]{R_2-2R_1}
\left(\begin{array}{ccc|c}
1&1&1&3\\
0&-3&-1&-4\\
0&-3&-1&-4
\end{array}\right)
\xrightarrow{R_3-R_2}
\left(\begin{array}{ccc|c}
1&1&1&3\\
0&-3&-1&-4\\
0&0&0&0
\end{array}\right).
\]
故 $r(A)=r(A\mid b)=2<3=n$，方程组有无穷多解。由第二行得
\[
-3x_2-x_3=-4\ \Longrightarrow\ x_3=4-3x_2,
\]
再由第一行得
\[
x_1=3-x_2-x_3=3-x_2-(4-3x_2)=2x_2-1.
\]
令自由未知量 $x_2=t$，得
\[
x=\begin{pmatrix}-1\\0\\4\end{pmatrix}+t\begin{pmatrix}2\\1\\-3\end{pmatrix},\qquad t\in\mathbb R.
\]
其中 $\eta^*=(-1,0,4)^T$ 是特解（$t=0$），$\xi=(2,1,-3)^T$ 是导出组的基础解系（$n-r=1$）。

代回验证：$2(-1+2t)-t+(4-3t)=2$，$3(-1+2t)+2(4-3t)=5$，均成立。

\boxed{\text{答案：}x=(-1,0,4)^T+t(2,1,-3)^T,\ t\in\mathbb R}
\else \answerarea[8] \fi%

\textbf{4.} 设 $Ax=b$ 是含有 $3$ 个未知数的非齐次线性方程组，$r(A)=1$，且
\[
\eta_1=(1,0,0)^T,\qquad \eta_2=(0,1,0)^T,\qquad \eta_3=(0,0,1)^T
\]
都是它的解。求 $Ax=b$ 的通解。

\ifshowsolutions%
\boxed{\text{解}}

\textbf{思路：}非齐次组的两个解之差是导出组的解；导出组的维数为 $n-r$，只需找到同样多个线性无关的差向量。

由 $A\eta_i=b$ 得
\[
A(\eta_2-\eta_1)=b-b=0,\qquad A(\eta_3-\eta_1)=0,
\]
即 $\eta_2-\eta_1=(-1,1,0)^T$ 与 $\eta_3-\eta_1=(-1,0,1)^T$ 都是导出组 $Ax=0$ 的解。

考察由它们组成的矩阵
\[
\begin{pmatrix}-1&-1\\1&0\\0&1\end{pmatrix},
\]
其前两行构成的行列式 $\begin{vmatrix}-1&-1\\1&0\end{vmatrix}=1\neq0$，故两向量线性无关。

又 $r(A)=1$，$n=3$，导出组解空间维数为 $n-r(A)=2$，所以
\[
\xi_1=(-1,1,0)^T,\qquad \xi_2=(-1,0,1)^T
\]
构成导出组的一个基础解系。于是
\[
x=\eta_1+k_1(\eta_2-\eta_1)+k_2(\eta_3-\eta_1)
=(1,0,0)^T+k_1(-1,1,0)^T+k_2(-1,0,1)^T.
\]
例如 $A=\begin{pmatrix}1&1&1\\0&0&0\\0&0&0\end{pmatrix}$，$b=(1,0,0)^T$ 就满足题设，可见这样的方程组确实存在。

\boxed{\text{答案：}x=(1,0,0)^T+k_1(-1,1,0)^T+k_2(-1,0,1)^T,\ k_1,k_2\in\mathbb R}
\else \answerarea[8] \fi%

\textbf{5.} 设 $A$ 为 $3\times 3$ 矩阵。已知增广矩阵经初等行变换化为下列行最简形，分别判断对应方程组的解的情况；有解时求出通解。
\[
\text{(i)}\ \left(\begin{array}{ccc|c}
1&0&2&1\\
0&1&-1&2\\
0&0&0&3
\end{array}\right),
\qquad
\text{(ii)}\ \left(\begin{array}{ccc|c}
1&0&2&1\\
0&1&-1&2\\
0&0&0&0
\end{array}\right).
\]

\ifshowsolutions%
\boxed{\text{解}}

\textbf{思路：}行最简形中若有形如 $(0,0,0\mid c)$（$c\neq0$）的矛盾行，则无解；否则 $r(A)=r(A\mid b)$，自由未知量个数为 $n-r$。

（i）第三行对应方程 $0x_1+0x_2+0x_3=3$，即 $0=3$，矛盾，故方程组\textbf{无解}。此时 $r(A)=2$，$r(A\mid b)=3$。

（ii）没有矛盾行，$r(A)=r(A\mid b)=2<3$，方程组有\textbf{无穷多解}。同解方程组为
\[
\begin{cases}
x_1+2x_3=1,\\
x_2-x_3=2.
\end{cases}
\]
取 $x_3=t$，得 $x_1=1-2t$，$x_2=2+t$，故
\[
x=(1,2,0)^T+t(-2,1,1)^T,\qquad t\in\mathbb R.
\]
导出组的基础解系为 $(-2,1,1)^T$（$n-r=1$）。

\boxed{\text{答案：(i) 无解；(ii) }x=(1,2,0)^T+t(-2,1,1)^T}
\else \answerarea[8] \fi%

\subsection{综合提高题}

\textbf{6.} 讨论参数 $a,b$ 取何值时，线性方程组
\[
\begin{cases}
x_1+x_2+x_3=1,\\
x_1+2x_2+3x_3=3,\\
x_1+3x_2+ax_3=b
\end{cases}
\]
无解、有唯一解、有无穷多解；有解时求出其解。

\ifshowsolutions%
\boxed{\text{解}}

\textbf{思路：}系数矩阵是方阵，先算 $\det A$ 找分界点；在 $\det A=0$ 的分界值处再化增广矩阵，用 $b$ 判断是否出现矛盾行。

系数行列式
\[
\det A=\begin{vmatrix}
1&1&1\\
1&2&3\\
1&3&a
\end{vmatrix}
=1\cdot(2a-9)-1\cdot(a-3)+1\cdot(3-2)=a-5.
\]

\textbf{情形一：$a\neq5$。}此时 $\det A\neq0$，$r(A)=r(A\mid b)=3=n$，方程组有唯一解。对增广矩阵消元：
\[
(A\mid b)\to
\left(\begin{array}{ccc|c}
1&1&1&1\\
0&1&2&2\\
0&2&a-1&b-1
\end{array}\right)
\xrightarrow{R_3-2R_2}
\left(\begin{array}{ccc|c}
1&1&1&1\\
0&1&2&2\\
0&0&a-5&b-5
\end{array}\right).
\]
回代得
\[
x_3=\frac{b-5}{a-5},\qquad
x_2=2-2x_3=\frac{2(a-b)}{a-5},\qquad
x_1=1-x_2-x_3=\frac{b-a}{a-5}.
\]

\textbf{情形二：$a=5$。}此时第三行化为 $(0,0,0\mid b-5)$。
若 $b=5$，则 $r(A)=r(A\mid b)=2<3$，方程组有无穷多解。同解方程组为
\[
x_1+x_2+x_3=1,\qquad x_2+2x_3=2.
\]
取 $x_3=t$，得 $x_2=2-2t$，$x_1=1-(2-2t)-t=t-1$，即
\[
x=(-1,2,0)^T+t(1,-2,1)^T,\qquad t\in\mathbb R.
\]
若 $b\neq5$，则 $r(A)=2$，$r(A\mid b)=3$，方程组无解。

\boxed{\begin{array}{l}
\text{当 }a\neq5\text{ 时，有唯一解 }x=\dfrac{1}{a-5}(b-a,\ 2(a-b),\ b-5)^{T};\\[4pt]
\text{当 }a=5,\ b=5\text{ 时，有无穷多解 }x=(-1,2,0)^{T}+t(1,-2,1)^{T};\\[4pt]
\text{当 }a=5,\ b\neq5\text{ 时，无解。}
\end{array}}
\else \answerarea[10] \fi%

\textbf{7.} 设 $A$ 为 $m\times n$ 矩阵，$r(A)=r<n$，非齐次方程组 $Ax=b$（$b\neq0$）有解，$\eta^*$ 是它的一个特解，$\xi_1,\xi_2,\ldots,\xi_{n-r}$ 是导出组 $Ax=0$ 的一个基础解系。证明：
（1）$\eta^*,\ \eta^*+\xi_1,\ \ldots,\ \eta^*+\xi_{n-r}$ 线性无关；
（2）$Ax=b$ 的任一解都可表示为这 $n-r+1$ 个向量的仿射组合（即组合系数之和为 $1$）；
（3）$Ax=b$ 的解集中任意 $n-r+2$ 个向量必线性相关，从而解集中线性无关的向量最多 $n-r+1$ 个。

\ifshowsolutions%
\boxed{\text{解}}

\textbf{思路：}证明线性无关时，先作线性组合，再用 $A$ 作用一次，把非齐次部分 $b$ 消掉，剩下的关系交给基础解系的线性无关性。

（1）设
\[
c_0\eta^*+c_1(\eta^*+\xi_1)+\cdots+c_{n-r}(\eta^*+\xi_{n-r})=0.
\]
整理得
\[
\Bigl(\sum_{i=0}^{n-r}c_i\Bigr)\eta^*+\sum_{i=1}^{n-r}c_i\xi_i=0.
\]
两边左乘 $A$，利用 $A\eta^*=b$、$A\xi_i=0$，得
\[
\Bigl(\sum_{i=0}^{n-r}c_i\Bigr)b=0.
\]
因 $b\neq0$，故 $\sum_{i=0}^{n-r}c_i=0$。代回上式得 $\sum_{i=1}^{n-r}c_i\xi_i=0$。由于 $\xi_1,\ldots,\xi_{n-r}$ 线性无关，$c_1=\cdots=c_{n-r}=0$，再由系数和为零得 $c_0=0$。因此这 $n-r+1$ 个向量线性无关。

（2）设 $\eta$ 是 $Ax=b$ 的任一解，则 $\eta-\eta^*$ 是导出组的解，故存在常数 $k_1,\ldots,k_{n-r}$ 使
\[
\eta=\eta^*+\sum_{i=1}^{n-r}k_i\xi_i.
\]
将其改写为
\[
\eta=\Bigl(1-\sum_{i=1}^{n-r}k_i\Bigr)\eta^*+\sum_{i=1}^{n-r}k_i(\eta^*+\xi_i),
\]
各系数之和为 $\bigl(1-\sum k_i\bigr)+\sum k_i=1$，即为仿射组合。

（3）设 $v_1,\ldots,v_{n-r+2}$ 是解集中的任意 $n-r+2$ 个向量。对 $j=2,\ldots,n-r+2$，有
\[
A(v_j-v_1)=b-b=0,
\]
即 $v_j-v_1\in\ker A$。而 $\dim\ker A=n-r$，故这 $n-r+1$ 个向量 $v_2-v_1,\ldots,v_{n-r+2}-v_1$ 线性相关，存在不全为零的 $c_2,\ldots,c_{n-r+2}$ 使
\[
\sum_{j=2}^{n-r+2}c_j(v_j-v_1)=0.
\]
展开得
\[
\sum_{j=2}^{n-r+2}c_jv_j-\Bigl(\sum_{j=2}^{n-r+2}c_j\Bigr)v_1=0.
\]
这是 $v_1,\ldots,v_{n-r+2}$ 的一个非平凡线性组合（系数 $c_2,\ldots$ 不全为零），故它们线性相关。因此解集中线性无关向量最多为 $n-r+1$ 个。

\boxed{\text{答案：上述（1）（2）（3）均得证；解集中线性无关向量最多 }n-r+1\text{ 个。}}
\else \answerarea[10] \fi%

\textbf{8.} 设齐次线性方程组
\[
\begin{cases}
x_1+x_2+x_3=0,\\
x_1+2x_2+ax_3=0,\\
x_1+4x_2+a^2x_3=0.
\end{cases}
\]
问 $a$ 取何值时方程组只有零解？$a$ 取何值时有非零解？有非零解时，求其基础解系与通解。

\ifshowsolutions%
\boxed{\text{解}}

\textbf{思路：}$3$ 个未知数的齐次方程组有非零解 $\iff$ 系数行列式为零。先算行列式定出 $a$ 的分界值，再逐个代入求基础解系。

系数行列式
\[
\det A=\begin{vmatrix}
1&1&1\\
1&2&a\\
1&4&a^2
\end{vmatrix}
\xrightarrow[R_3-R_1]{R_2-R_1}
\begin{vmatrix}
1&1&1\\
0&1&a-1\\
0&3&a^2-1
\end{vmatrix}
=(a^2-1)-3(a-1)=(a-1)(a-2).
\]

\textbf{（1）$a\neq1$ 且 $a\neq2$。}$\det A\neq0$，即 $r(A)=3=n$，方程组只有零解。

\textbf{（2）$a=1$。}此时
\[
A=\begin{pmatrix}1&1&1\\1&2&1\\1&4&1\end{pmatrix}
\xrightarrow[R_3-R_1]{R_2-R_1}
\begin{pmatrix}1&1&1\\0&1&0\\0&3&0\end{pmatrix}
\to
\begin{pmatrix}1&0&1\\0&1&0\\0&0&0\end{pmatrix}.
\]
$r(A)=2$，$n-r=1$。由 $x_2=0$，$x_1+x_3=0$，取 $x_3=1$ 得基础解系 $\xi=(1,0,-1)^T$，通解 $x=k(1,0,-1)^T$。

\textbf{（3）$a=2$。}此时
\[
A=\begin{pmatrix}1&1&1\\1&2&2\\1&4&4\end{pmatrix}
\xrightarrow[R_3-R_1]{R_2-R_1}
\begin{pmatrix}1&1&1\\0&1&1\\0&3&3\end{pmatrix}
\to
\begin{pmatrix}1&0&0\\0&1&1\\0&0&0\end{pmatrix}.
\]
$r(A)=2$，$n-r=1$。由 $x_1=0$，$x_2+x_3=0$，取 $x_3=1$ 得基础解系 $\xi=(0,-1,1)^T$，通解 $x=k(0,-1,1)^T$。

\boxed{\text{答案：}a\neq1,2\text{ 时仅有零解；}a=1\text{ 时 }x=k(1,0,-1)^{T};\ a=2\text{ 时 }x=k(0,-1,1)^{T}}
\else \answerarea[10] \fi%

\textbf{9.} 设 $A$ 为 $m\times n$ 矩阵。证明：对任意 $m$ 维列向量 $b$，方程组 $Ax=b$ 都有解的充要条件是 $r(A)=m$。并回答：
（1）此时 $m$ 与 $n$ 有什么关系？举例说明 $m<n$ 时确实可能对每个 $b$ 都有解；
（2）若对任意 $b$，$Ax=b$ 都有唯一解，则 $A$ 是什么形状的矩阵？

\ifshowsolutions%
\boxed{\text{解}}

\textbf{思路：}把“对每个 $b$ 都有解”翻译成“$A$ 的列空间等于整个 $\mathbb R^m$”，再用列空间维数等于 $r(A)$ 建立等式。

$Ax=b$ 有解等价于 $b$ 可由 $A$ 的列向量线性表出，即 $b\in\operatorname{Col}(A)$。因此“对任意 $b\in\mathbb R^m$ 都有解”等价于
\[
\operatorname{Col}(A)=\mathbb R^m.
\]
而 $\dim\operatorname{Col}(A)=r(A)$，故
\[
\operatorname{Col}(A)=\mathbb R^m
\iff \dim\operatorname{Col}(A)=m
\iff r(A)=m.
\]
必要性：若对每个 $b$ 都有解，则 $\mathbb R^m=\operatorname{Col}(A)$，取维数得 $r(A)=m$。充分性：若 $r(A)=m$，则 $\operatorname{Col}(A)$ 是 $\mathbb R^m$ 的 $m$ 维子空间，只能是 $\mathbb R^m$ 本身，故每个 $b$ 都有解。

（1）由 $r(A)\le n$ 得 $m=r(A)\le n$。当 $m<n$ 时确实可能成立，例如
\[
A=\begin{pmatrix}1&0&0\\0&1&0\end{pmatrix},
\]
$m=2$，$n=3$，$r(A)=2=m$，对任意 $b=(b_1,b_2)^T$，取 $x=(b_1,b_2,0)^T$ 即可。

（2）若每个 $b$ 都有唯一解，则一方面由上面的结论 $r(A)=m$；另一方面唯一性要求导出组只有零解，即 $n-r(A)=0$，故 $r(A)=n$。于是 $m=n=r(A)$，$A$ 为 $n$ 阶方阵且满秩，即 $A$ 可逆。

\boxed{\text{答案：}r(A)=m\text{；此时 }m\le n\text{；若解还唯一，则 }m=n\text{ 且 }A\text{ 可逆}}
\else \answerarea[10] \fi%

\subsection{真题风格与综合训练}

\textbf{10.} \textbf{（选择题）}设 $A$ 为 $3\times 4$ 矩阵，且 $r(A)=2$，则（\quad）
\begin{center}
(A) 对任意 $b\in\mathbb R^3$，$Ax=b$ 有解 \\[3pt]
(B) 存在 $b\in\mathbb R^3$，使 $Ax=b$ 有唯一解 \\[3pt]
(C) 若 $Ax=b$ 有解，则必有无穷多解 \\[3pt]
(D) $Ax=0$ 只有零解
\end{center}

\ifshowsolutions%
\boxed{\text{解}}

\textbf{思路：}用秩判定逐项排查：有解看 $r(A)$ 与 $r(A\mid b)$ 是否相等，解唯一性看 $r(A)$ 是否等于未知数个数 $4$。

$n=4$，$r(A)=2$，导出组解空间维数为 $n-r(A)=2$。

(A) 错。对任意 $b$ 都有解要求 $r(A)=m=3$，而这里 $r(A)=2<3$，列空间只是 $\mathbb R^3$ 的二维子空间，必存在不在其中的 $b$。

(B) 错。唯一解要求 $r(A)=n=4$，与 $r(A)=2$ 矛盾。

(C) 对。若 $Ax=b$ 有解，则 $r(A\mid b)=r(A)=2<4=n$，自由未知量有 $2$ 个，故有无穷多解。

(D) 错。$n-r(A)=2>0$，导出组有非零解。

\boxed{\text{答案：}(C)}
\else \answerarea[4] \fi%

\textbf{11.} \textbf{（填空题）}设 $A$ 为 $m\times 4$ 矩阵，$r(A)=3$，且
\[
\eta_1=(1,1,0,0)^T,\qquad \eta_2=(0,1,1,0)^T
\]
都是 $Ax=b$ 的解，则 $Ax=b$ 的通解为 \underline{\hspace{5cm}}，导出组 $Ax=0$ 的基础解系为 \underline{\hspace{4cm}}。

\ifshowsolutions%
\boxed{\text{解}}

\textbf{思路：}非齐次组的两个解之差是导出组的解；导出组的基础解系应含 $n-r=4-3=1$ 个向量，故只需一个非零的差向量。

由 $A\eta_1=A\eta_2=b$ 得
\[
A(\eta_2-\eta_1)=b-b=0,
\]
即 $\eta_2-\eta_1=(-1,0,1,0)^T$ 是导出组的解。它显然非零，而导出组解空间维数为 $4-3=1$，所以它就是导出组的一个基础解系。于是非齐次组的通解为
\[
x=\eta_1+k(\eta_2-\eta_1)=(1,1,0,0)^T+k(-1,0,1,0)^T,\qquad k\in\mathbb R.
\]

\boxed{\text{答案：}x=(1,1,0,0)^{T}+k(-1,0,1,0)^{T};\ \text{基础解系为 }\{(-1,0,1,0)^{T}\}}
\else \answerarea[3] \fi%

\textbf{12.} \textbf{（解答题）}已知线性方程组
\[
\begin{cases}
\lambda x_1+x_2+x_3=1,\\
x_1+\lambda x_2+x_3=1,\\
x_1+x_2+\lambda x_3=1.
\end{cases}
\]
讨论 $\lambda$ 取何值时方程组无解、有唯一解、有无穷多解，并在有无穷多解时求出通解。

\ifshowsolutions%
\boxed{\text{解}}

\textbf{思路：}系数矩阵的主对角线全为 $\lambda$、其余元素全为 $1$，其特征值为 $\lambda+2$（一重）与 $\lambda-1$（二重），据此写出行列式并定出分界点 $\lambda=1,-2$，再逐值代入。

系数行列式
\[
\det A=(\lambda+2)(\lambda-1)^2.
\]
直接按第一行展开即可得到该分解。

\textbf{（1）$\lambda\neq1$ 且 $\lambda\neq-2$。}$\det A\neq0$，$r(A)=r(A\mid b)=3=n$，方程组有唯一解。由于方程组在交换任意两个未知量（同时交换对应方程）后不变，由唯一性知 $x_1=x_2=x_3$。代入第一式得
\[
(\lambda+2)x_1=1\ \Longrightarrow\ x_1=x_2=x_3=\frac{1}{\lambda+2}.
\]

\textbf{（2）$\lambda=1$。}三个方程都化为
\[
x_1+x_2+x_3=1.
\]
$r(A)=r(A\mid b)=1<3$，方程组有无穷多解。取 $x_2=s,\ x_3=t$，则 $x_1=1-s-t$，通解为
\[
x=(1,0,0)^T+s(-1,1,0)^T+t(-1,0,1)^T,\qquad s,t\in\mathbb R.
\]
导出组的基础解系为 $(-1,1,0)^T,(-1,0,1)^T$。

\textbf{（3）$\lambda=-2$。}此时 $r(A)=2$（例如前两行不成比例）。将三个方程相加，左边为
\[
(-2+1+1)(x_1+x_2+x_3)=0,
\]
右边为 $3$，得到矛盾 $0=3$，故 $r(A\mid b)=3\neq r(A)$，方程组无解。

\boxed{\begin{array}{l}
\text{当 }\lambda\neq1,-2\text{ 时，有唯一解 }x_1=x_2=x_3=\dfrac{1}{\lambda+2};\\[4pt]
\text{当 }\lambda=1\text{ 时，有无穷多解 }x=(1,0,0)^{T}+s(-1,1,0)^{T}+t(-1,0,1)^{T};\\[4pt]
\text{当 }\lambda=-2\text{ 时，无解。}
\end{array}}
\else \answerarea[10] \fi%
